% Generated by code/assemble.py; edit part_B_audit.tex.
\subsection{Arithmetic of the varying family}

\begin{result}{A.101}{Refuted}{Varying ghost families are not compact; correction to Lead L2}
The original A.44 family is noncompact in $\mathcal K^2$. In particular
the set of ordinary ghosts surviving the integral class-one quotient
is noncompact. This refutes the compactness assertion in Lead L2;
it does not refute compactness of section fibres over a fixed abelian
cocycle. More generally, if the fixed-depth existence assertion of
GT\_TOWER\_CHECK C.2.12 is retained, the set of ordinary ghosts
surviving any prescribed finite integral class is noncompact.
\end{result}
\begin{proof}
For the original family, reduction modulo fifth powers gives
\[
A_5=[3q_0^{24}]_5,
\]
because $5\mid(j+4)!$ for every $j\ge1$. E.102 permits infinitely
many different $q_0$. Their valuation vectors modulo five are
different, since $24=4\pmod5$ and $q_0>3$. Thus the image in the
discrete fifth-power quotient is infinite. A compact space has finite
image in a discrete space, giving the contradiction. At class one
every pair is the abelian cocycle of a section relative to the
tangential splitting, so this example lies in the asserted survivor
set.

For the last assertion, fix a class $r$ and apply C.2.12 with
$(R,N)=(r,k!)$. It produces ghosts in the class-$r$ survivor set whose
abelian coordinates agree with $(3,-2)$ modulo $k!$. They converge in
$\mathcal K^2$ to that rational pair, which is not a ghost. Hence this
survivor set is not closed in a Hausdorff space, and cannot be compact.
This second argument is explicitly conditional on C.2.12; the first
argument requires only the arithmetic construction E.102.

The fixed-pair repair is different in kind: if the literal continuous
section fibres above one fixed abelian cocycle are compact and the
transition maps are continuous, nonemptiness at each class implies
nonemptiness of their inverse limit by the finite-intersection
property. Nothing in this argument changes, or chooses, the abelian
pair. It supplies no compactness theorem for the varying ghost locus.
\end{proof}
